% ================================================================
%  PDL --- D39
%  Derivation of kappa = R_surf / R_tot from the PDL Axioms:
%  The Indifference Lemma and Partial Resolution of OP1 (D36)
%  C. Laubscher, 2026
% ================================================================
\documentclass[11pt,a4paper]{article}
\usepackage{mathtools}
\usepackage[utf8]{inputenc}
\usepackage[T1]{fontenc}
\usepackage[british]{babel}
\usepackage{amsmath,amssymb,amsthm}
\usepackage{mathrsfs}
\usepackage[colorlinks=true,citecolor=blue,linkcolor=blue,urlcolor=blue]{hyperref}
\usepackage[margin=2.5cm]{geometry}
\usepackage{booktabs}
\usepackage{enumitem}
\usepackage{array}
\setlist[itemize,1]{label=---}
\usepackage{csquotes}
\usepackage{xcolor}

% ── Theorem environments ──────────────────────────────────────────────────────
\newtheorem{theorem}{Theorem}[section]
\newtheorem{proposition}[theorem]{Proposition}
\newtheorem{lemma}[theorem]{Lemma}
\newtheorem{corollary}[theorem]{Corollary}
\theoremstyle{definition}
\newtheorem{definition}[theorem]{Definition}
\newtheorem{remark}[theorem]{Remark}
\newtheorem{hypothesis}[theorem]{Hypothesis}
\newtheorem{openprob}{Open Problem}

% ── PDL macros ────────────────────────────────────────────────────────────────
\newcommand{\PDL}{\textnormal{PDL}}
\newcommand{\Kfour}{K_4}
\newcommand{\rval}{r_{\mathrm{val}}}
\newcommand{\Rsea}{R_{\mathrm{sea}}}
\newcommand{\Rtot}{R_{\mathrm{tot}}}
\newcommand{\Rsurf}{R_{\mathrm{surf}}}
\newcommand{\Gval}{G_{\mathrm{val}}}
\newcommand{\Gsea}{G_{\mathrm{sea}}}
\newcommand{\vp}{\varphi}
\newcommand{\kap}{\kappa}
\newcommand{\sig}{\sigma}
\newcommand{\GPDL}{G_{\mathrm{PDL}}}
\newcommand{\Geff}{G_{\mathrm{eff}}}
\newcommand{\hb}{\hbar}

% ================================================================
\begin{document}

\title{\textbf{Derivation of $\kappa = R_{\mathrm{surf}}/R_{\mathrm{tot}}$\\
from the \PDL\ Axioms:\\
The Indifference Lemma and Partial Resolution\\
of Open Problem~OP1 (D36)}\\[4pt]
\large Document D39 --- PDL Programme}
\author{C\'edric Laubscher\\
\small Independent Researcher, Switzerland}
\date{March 2026}

\maketitle

% ── Abstract ──────────────────────────────────────────────────────────────────
\begin{abstract}
Open Problem~OP1 of D36 asks for a derivation of $\kap = \Rsurf/\Rtot$
from the four \PDL\ axioms~C1--C4, without invoking D25 as a definition.
The present paper provides a systematic investigation of this problem and
reaches the following conclusions.

First, a \emph{two-factor decomposition} is established:
$\kap = (\vp/3)\cdot(\rval/\Rtot)$, where the first factor
$P_1 = \vp/3 = \Rsurf/\rval$ is the conditional engagement probability
of a valence relation (established in~D05), and the second factor
$P_2 = \rval/\Rtot$ is the probability that a relation drawn uniformly
from the full relational budget is a valence relation. The identity
$\kap = P_1 \cdot P_2$ is verified to machine precision
(error~$< 10^{-15}$).

Second, a \emph{negative result} is proved: axiom~C3 alone does not
imply the uniform measure on $\Rtot$. A constructive counterexample
(a graph satisfying~C3 with a structurally asymmetric edge) demonstrates
this explicitly.

Third, a \emph{positive result under a named hypothesis} is established.
Under Hypothesis~H3 (the Indifference Lemma: the measure on the
$\Rtot$ relations of a \PDL\ closure is uniform, as a consequence
of the absence of any external distinction between relations under
axioms~C3 and~C4), the full decomposition $\kap = P_1 \cdot P_2$
follows from the \PDL\ axioms together with D05 and D29.
Hypothesis~H3 is shown to be: (i)~consistent with all \PDL\ axioms;
(ii)~the unique measure compatible with the $S_4$-symmetry of $\Kfour$
(proved in~D36); and (iii)~constrained to better than $0.01\%$ by the
Hubble tension observable.

OP1~D36 is therefore \emph{partially resolved}: the derivation is
complete under H3, and H3 is the unique residual hypothesis.
\end{abstract}

\newpage

\tableofcontents
\newpage

% ══════════════════════════════════════════════════════════════════════════════
\section{Introduction and position in the programme}
\label{sec:intro}

\subsection{The open problem}

Paper~D36 of the \PDL\ corpus established the Gate~3 theorem
$\Geff(N) = \sig(N)\cdot\GPDL$ under two named hypotheses: H1
($\kap = \Rsurf/\Rtot$ as gravitational engagement fraction) and H2
(proportionality of $G$ to engagement fraction). It identified OP1~D36
as the derivation of H1 from axioms~C1--C4 alone.

Concretely, H1 states:
\begin{quote}
$\kap = \Rsurf/\Rtot$ is the probability that a relation drawn uniformly
from the $\Rtot$ relations of a \PDL\ proton closure is gravitationally
active under the $(A)\wedge(B)$ stability criterion.
\end{quote}

The present paper investigates this statement systematically. It
establishes what can and cannot be proved from the axioms, names the
residual hypothesis precisely, and demonstrates that this residual is
the unique bottleneck.

\subsection{The two-factor structure of \texorpdfstring{$\kap$}{kappa}}

The key observation, established numerically in Section~\ref{sec:decomp}
and analytically in Section~\ref{sec:proof}, is:

\begin{equation}
\kap = \underbrace{\frac{\Rsurf}{\rval}}_{P_1 = \vp/3}
\;\times\;
\underbrace{\frac{\rval}{\Rtot}}_{P_2},
\label{eq:decomp_intro}
\end{equation}

where:
\begin{itemize}
\item $P_1 = \Rsurf/\rval = \vp/3$ is the probability that a valence
  relation is gravitationally active. This was established in D05 via
  the golden-ratio redistribution of coherence, and is a consequence
  of axiom~C4.
\item $P_2 = \rval/\Rtot$ is the probability that a relation drawn
  uniformly from $\Rtot$ is a valence relation. This requires the
  uniform measure on $\Rtot$, which is Hypothesis~H3.
\end{itemize}

The paper is organised as follows.
Section~\ref{sec:prereq} recalls the prerequisites.
Section~\ref{sec:decomp} establishes the two-factor decomposition
numerically and analytically.
Section~\ref{sec:negative} proves the negative result (C3 alone
is insufficient).
Section~\ref{sec:H3} formulates H3, proves its consistency and
uniqueness, and establishes the main theorem.
Section~\ref{sec:constraint} quantifies the observational constraint
on H3.
Section~\ref{sec:epistemic} gives the epistemic status table.

% ══════════════════════════════════════════════════════════════════════════════
\section{Prerequisites}
\label{sec:prereq}

\subsection{The four PDL axioms}

The four axioms of the \PDL\ framework~\cite{PDL_Main_2026} are:
\begin{description}
\item[Axiom~1 (Binary pulsation).] Existence is a reproducible
  two-state alternation. There exists a sign configuration $s^*$
  and a map $\Phi$ such that $\Phi(\Phi(s^*)) = s^* \neq \Phi(s^*)$.
\item[Axiom~2 (Triangular coherence).] Stable configurations minimise
  the fraction $\nu$ of violated triangles. A triangle $(i,j,k)$ is
  coherent iff $s_{ij}s_{jk}s_{ik} = +1$.
\item[Axiom~3 (Minimal completeness / closure, C3).] All relational
  references necessary to sustain persistence are internal to the
  structure. No external background is presupposed. Formally: for
  every vertex $v$, there exists at least one coherent triangle
  containing $v$, and the structure is self-referentially closed.
\item[Axiom~4 (Logical optimisation, C4).] Among admissible
  configurations, only those maximising the optimisation functional
  $\Omega(\sigma) = w_\mathrm{coh}(1-\nu) + w_\mathrm{conn}\rho
  + w_\mathrm{min}\delta_\mathrm{min}$ are stable.
\end{description}

\subsection{The proton quintuplet and relational budget}

The proton \PDL\ closure is characterised by
$(n_u, n_d, \rval, \Rsea, \Rtot) = (24, 28, 930, 10087, 11017)$,
with $\rval = 2\binom{n_u}{2} + \binom{n_d}{2}$ and
$\Rtot = \rval + \Rsea$~\cite{PDL_Main_2026,Combinatorial_Proton_PDL_2026}.
The active surface is $\Rsurf = 310\vp \in \mathbb{Q}(\sqrt{5})$
(D05~\cite{PDL_Main_2026}).

\subsection{the \texorpdfstring{$(A)\wedge(B)$}{(A) wedge (B)} criterion and its consequences}

D29~\cite{PDL_D29_2026} established:
\begin{itemize}
\item A mixed triangle $(e_i, e_j, p_k)$ is stable iff
  $(A)\wedge(B)$: $P_1 = s_{K_4}^{(1)}(e_i,e_j)$ and
  $P_2 = -P_1$.
\item $\Gval$ satisfies $(A)\wedge(B)$ structurally: exactly
  $\Rsurf$ of the $\rval$ valence relations are stably active.
\item $\Gsea$ satisfies $(A)\wedge(B)$ with probability $(1/4)^N \to 0$
  in the stationary regime.
\end{itemize}

These results hold without any assumption about the measure on $\Rtot$.

% ══════════════════════════════════════════════════════════════════════════════
\section{The two-factor decomposition of \texorpdfstring{$\kap$}{kappa}}
\label{sec:decomp}

\subsection{Analytical derivation}

\begin{proposition}[Two-factor decomposition]
\label{prop:decomp}
The engagement fraction $\kap = \Rsurf/\Rtot$ decomposes as:
\begin{equation}
\kap = P_1 \cdot P_2,
\qquad
P_1 = \frac{\Rsurf}{\rval} = \frac{\vp}{3},
\quad
P_2 = \frac{\rval}{\Rtot}.
\label{eq:decomp}
\end{equation}
\end{proposition}

\begin{proof}
Direct computation:
$P_1 \cdot P_2 = (\Rsurf/\rval)\cdot(\rval/\Rtot) = \Rsurf/\Rtot = \kap$.
\end{proof}

\begin{remark}[Physical interpretation]
The decomposition~\eqref{eq:decomp} admits a natural probabilistic
reading under the hypothesis of a uniform measure on $\Rtot$
(Hypothesis~H3, Section~\ref{sec:H3}):
\begin{align}
P_2 &= P(\text{relation is a valence relation})
= \frac{\rval}{\Rtot}, \label{eq:P2} \\
P_1 &= P(\text{relation is stably active} \mid
\text{relation is a valence relation})
= \frac{\Rsurf}{\rval} = \frac{\vp}{3}. \label{eq:P1}
\end{align}
By the law of total probability, and using
$P(\text{active} \mid \text{sea}) = 0$ (D29):
\begin{equation}
\kap = P(\text{active}) = P_1 \cdot P_2
+ \underbrace{0}_{\text{sea}} \cdot \frac{\Rsea}{\Rtot}
= \frac{\Rsurf}{\Rtot}.
\label{eq:total_prob}
\end{equation}
\end{remark}

\subsection{Numerical verification}

Table~\ref{tab:decomp} gives the numerical values of all quantities
and verifies equation~\eqref{eq:decomp} to machine precision.

\begin{table}[ht]
\centering
\caption{Numerical verification of the two-factor decomposition.}
\label{tab:decomp}
\medskip
\begin{tabular}{lrr}
\toprule
Quantity & Formula & Value \\
\midrule
$\vp$ & $(1+\sqrt{5})/2$ & $1.61803399$ \\
$\Rsurf$ & $310\vp$ & $501.59054$ \\
$\rval$ & $930$ & $930$ \\
$\Rtot$ & $11017$ & $11017$ \\
$P_1 = \Rsurf/\rval$ & $\vp/3$ & $0.53934466$ \\
$P_2 = \rval/\Rtot$ & $930/11017$ & $0.08441500$ \\
$P_1 \cdot P_2$ & & $0.04552878$ \\
$\kap = \Rsurf/\Rtot$ & $310\vp/11017$ & $0.04552878$ \\
Error & $|P_1 P_2 - \kap|$ & $< 10^{-15}$ \\
\bottomrule
\end{tabular}
\end{table}

% ══════════════════════════════════════════════════════════════════════════════
\section{Negative result: C3 alone is insufficient}
\label{sec:negative}

\begin{proposition}[C3 does not imply uniform measure]
\label{prop:negative}
Axiom~C3 alone does not imply that the natural measure on the
relations of a \PDL\ closure is uniform on $\Rtot$.
\end{proposition}

\begin{proof}
We exhibit a constructive counterexample. Consider the signed graph
$G$ on 6 vertices with edge set:
\[
E = \{(0,1),(0,2),(1,2),(3,4),(3,5),(4,5),(2,3)\}.
\]
This graph satisfies C3: every vertex participates in at least one
coherent triangle (triangles $\{0,1,2\}$ and $\{3,4,5\}$ cover all
6 vertices). However, the edge $(2,3)$ participates in \emph{zero}
triangles, while all other edges participate in exactly one triangle.
The edge $(2,3)$ is therefore structurally asymmetric with respect
to the others: it cannot be assigned the same measure as the triangle
edges without violating the structural distinction imposed by C3 itself.

A uniform measure on $E$ would assign equal probability $1/7$ to each
edge. But the edge $(2,3)$ is the unique edge with zero triangle
participation, distinguishing it from all others. C3 does not eliminate
this asymmetry; it generates it. Therefore, C3 alone does not force
uniform measure on the full edge set. \qed
\end{proof}

\begin{remark}[Why this matters]
Proposition~\ref{prop:negative} shows that the derivation of
$\kap = \Rsurf/\Rtot$ cannot be completed from C3 alone. An additional
structure is required. This is provided by Hypothesis~H3, formulated
in Section~\ref{sec:H3}.
\end{remark}

% ══════════════════════════════════════════════════════════════════════════════
\section{Hypothesis H3: the Indifference Lemma}
\label{sec:H3}

\subsection{Formulation}

\begin{hypothesis}[H3 --- Indifference Lemma]
\label{hyp:H3}
Let $\mathcal{P}$ be a \PDL\ proton closure satisfying axioms~C1--C4,
with total relational budget $\Rtot$. The natural measure on the
$\Rtot$ relations of $\mathcal{P}$, as seen by an external
$\Kfour$ observer applying the $(A)\wedge(B)$ criterion, is the
uniform measure: each relation has measure $1/\Rtot$.
\end{hypothesis}

\subsection{Motivation from C3 and C4}

Although C3 alone does not imply H3 (Proposition~\ref{prop:negative}),
the combination of C3 and C4 provides strong structural motivation.

\begin{lemma}[C3 eliminates external distinctions]
\label{lem:C3}
Under axiom~C3, no relation of a \PDL\ closure has a privileged
status \emph{with respect to an external observer}. The internal
distinction between valence and sea relations is produced by C4
(optimisation), not by C3. An external $\Kfour$ block applying
$(A)\wedge(B)$ cannot distinguish valence from sea relations prior
to filtration: the criterion acts on the cross-edge signs of
$\Kfour$, which are independent of the internal valence/sea
partition.
\end{lemma}

\begin{proof}
The $(A)\wedge(B)$ criterion (D29) is:
$(A)$: $P_1 = s_{K_4}^{(1)}(e_i,e_j)$;
$(B)$: $P_2 = -P_1$,
where $P_c = s_\times(e_i,p_k,c)\cdot s_\times(e_j,p_k,c)$ is the
cross-product of $\Kfour$ signs with proton signs.
The cross-product depends on the signs of cross-edges between $\Kfour$
and the proton, not on whether the proton relation $p_k$ belongs to
$\Gval$ or $\Gsea$. The criterion is therefore blind to the
valence/sea partition. Under C3, which requires internal closure
without external reference, all $\Rtot$ relations are candidates
for cross-edge formation with equal standing from the perspective
of $\Kfour$.
\end{proof}

\begin{lemma}[C4 produces the valence/sea partition internally]
\label{lem:C4}
The distinction between $\Gval$ (stationary, $\rval = 930$ relations)
and $\Gsea$ (dynamical, $\Rsea = 10087$ relations) is produced by
axiom~C4 (optimisation of $\Omega$). This distinction is
\emph{internal} to the proton closure and is not visible to an
external $\Kfour$ observer before the $(A)\wedge(B)$ filtration is
applied.
\end{lemma}

\begin{proof}
The partition $\Rtot = \rval + \Rsea$ arises from the optimisation
of $\Omega(\sigma)$ under C4: stationary relations correspond to
local optima of $\Omega$, while dynamical relations sustain the
collective regime of the sea. This distinction is an output of C4,
not an input to the measurement process of an external $\Kfour$
observer. The $\Kfour$ block interacts with the proton via
cross-edges; the nature (valence or sea) of the proton relation
at the other end of a cross-edge is not accessible to $\Kfour$
until $(A)\wedge(B)$ is evaluated.
\end{proof}

\begin{remark}[Why C3 $\wedge$ C4 motivates but does not prove H3]
Lemmas~\ref{lem:C3} and~\ref{lem:C4} together establish that,
from the perspective of an external $\Kfour$ observer:
(i)~all $\Rtot$ relations are candidates for engagement (C3);
(ii)~no a priori distinction between them is available before
filtration (C4 makes the distinction internal).
This motivates H3 as a \emph{principle of indifference}: in the
absence of any distinguishing structure accessible to the external
observer, the natural measure is uniform.

However, this falls short of a strict algebraic proof: the argument
requires the identification of ``no external distinction available
to $\Kfour$'' with ``uniform measure'', which is the classical
principle of indifference. This principle is not itself contained
in C1--C4 as a formal axiom. H3 therefore remains a named hypothesis.
\end{remark}

\subsection{Consistency of H3}

\begin{proposition}[H3 is consistent with C1--C4]
\label{prop:consistent}
Hypothesis~H3 (uniform measure on $\Rtot$) is consistent with all
four \PDL\ axioms~C1--C4.
\end{proposition}

\begin{proof}
H3 makes a statement about the measure on the relations of a
\PDL\ closure as seen by an external observer. It does not constrain
the internal sign structure, the pulsation dynamics (C1--C2), the
closure property (C3), or the optimisation (C4). It adds a
probabilistic interpretation to an already well-defined combinatorial
structure, without contradicting any of the axioms.
\end{proof}

\subsection{Uniqueness of H3}

\begin{proposition}[H3 is the unique measure consistent with the
$S_4$ symmetry of $\Kfour$]
\label{prop:unique}
The uniform measure on $\Rtot$ is the unique $S_4$-invariant measure
on the relations of a \PDL\ proton closure, where $S_4 =
\mathrm{Aut}(\Kfour)$ acts on the proton via cross-edge permutations.
\end{proposition}

\begin{proof}
D36~\cite{PDL_D36_2026} proved that $S_4$ acts transitively on the
6 edges of $\Kfour$, with all 24 permutations verified exhaustively.
An $S_4$-invariant measure on the proton relations must assign
equal weight to any two relations that are related by an element
of $S_4$. Since $S_4$ acts transitively on the cross-edge connections
between $\Kfour$ and the proton, no two relations can be distinguished
by $S_4$-invariant quantities. The unique $S_4$-invariant measure is
therefore the uniform measure $1/\Rtot$ on each relation. Any
non-uniform measure would define a preferred direction in the
$\Kfour$ edge space, breaking $S_4$-invariance and contradicting
the transitivity proved in D36.
\end{proof}

\begin{remark}
Proposition~\ref{prop:unique} shows that H3 is not arbitrary: it is
the unique measure consistent with the symmetry of $\Kfour$ established
in D36. The residual openness of H3 is not that it could be replaced
by another measure, but that the equivalence between
``$S_4$-invariance of the coupling'' and ``uniform measure on $\Rtot$''
has not been derived from C1--C4 without invoking D36 as an additional
input.
\end{remark}

% ══════════════════════════════════════════════════════════════════════════════
\section{Main theorem: \texorpdfstring{$\kap$}{kappa} from H3, D05, D29}
\label{sec:proof}

\begin{theorem}[Derivation of $\kap$ under H3]
\label{thm:kappa}
Under Hypothesis~H3 (uniform measure on $\Rtot$), and given:
\begin{itemize}
\item D05~\cite{PDL_Main_2026}: $\Rsurf = \vp\,\rval/3$
  (golden-ratio redistribution under C4);
\item D29~\cite{PDL_D29_2026}: $P(\mathrm{stable}\mid\mathrm{sea}) = 0$
  in the stationary regime;
\end{itemize}
the gravitational engagement fraction of a single \PDL\ proton is:
\begin{equation}
\kap = \frac{\Rsurf}{\Rtot} = \frac{310\vp}{11017} \in \mathbb{Q}(\sqrt{5}).
\label{eq:kappa_theorem}
\end{equation}
\end{theorem}

\begin{proof}
Under H3, the probability that a relation drawn uniformly from $\Rtot$
is a valence relation is $P_2 = \rval/\Rtot$.
By D05, the conditional probability that a valence relation is stably
active under $(A)\wedge(B)$ is $P_1 = \Rsurf/\rval = \vp/3$.
By D29, the conditional probability that a sea relation is stably
active is $P(\mathrm{stable}\mid\mathrm{sea}) = 0$.
By the law of total probability:
\begin{align}
\kap &= P(\mathrm{stable})
= P(\mathrm{stable}\mid\mathrm{val})\cdot P(\mathrm{val})
+ P(\mathrm{stable}\mid\mathrm{sea})\cdot P(\mathrm{sea}) \notag\\
&= \frac{\Rsurf}{\rval}\cdot\frac{\rval}{\Rtot}
+ 0\cdot\frac{\Rsea}{\Rtot}
= \frac{\Rsurf}{\Rtot}
= \frac{310\vp}{11017}.
\label{eq:kappa_proof}
\end{align}
The result is in $\mathbb{Q}(\sqrt{5})$ since $\vp = (1+\sqrt{5})/2$
and all other factors are rational. \qed
\end{proof}

\begin{corollary}[Gate~3 under H1=H3, H2]
\label{cor:gate3}
Under H3 (which subsumes H1 of D36), H2, and the results of D31/D36,
the Gate~3 theorem holds:
$\Geff(N) = \sig(N)\cdot\GPDL$ with $\sig(N) = 1-(1-\kap)^N$
and $\kap = 310\vp/11017$.
\end{corollary}

% ══════════════════════════════════════════════════════════════════════════════
\section{Observational constraint on H3}
\label{sec:constraint}

H3 can be tested indirectly through its consequence $\kap = \Rsurf/\Rtot$.
Any deviation from H3 would produce a modified engagement fraction
$\kap' \neq \kap$, which would shift the Hubble tension prediction.

\begin{proposition}[H3 is constrained to $\pm 0.01\%$]
\label{prop:constraint}
The Hubble tension observable $H_{0,\mathrm{loc}}/H_{0,\mathrm{CMB}}
= 73.04/67.26 = 1.085935$ constrains $\kap$ to within $\pm 0.0001$
of $\kap_0 = 310\vp/11017$. This implies that any non-uniform
measure on $\Rtot$ is constrained to deviate from uniformity by less
than $0.22\%$ in the effective $\kap$.
\end{proposition}

\begin{proof}
A sensitivity analysis (D36, Table~2) shows that a deviation
$\delta\kap = 0.0001$ shifts the Hubble ratio by $0.030\%$, which
is $5\times$ the current $0.006\%$ agreement. Requiring the Hubble
ratio to remain within $2\sigma$ of the observed value constrains
$|\delta\kap| < 0.0001$. Since $\kap = \Rsurf/\Rtot = 0.04553$,
this is a relative constraint of $0.22\%$ on $\kap$, and equivalently
on the deviation from uniformity of the measure on $\Rtot$.
\end{proof}

Table~\ref{tab:sensitivity} summarises the sensitivity of the Hubble
ratio to deviations in $\kap$.

\begin{table}[ht]
\centering
\caption{Sensitivity of the Hubble ratio prediction to deviations
$\delta\kap$ from the PDL value $\kap_0 = 310\vp/11017$.
Data from D36~\cite{PDL_D36_2026}.}
\label{tab:sensitivity}
\medskip
\begin{tabular}{rrrrr}
\toprule
$\delta\kap$ & $\kap_0 + \delta\kap$ & Hubble ratio &
Residual from obs. & Status \\
\midrule
$0$ & $0.045529$ & $1.085868$ & $0.006\%$ & \textbf{PDL} \\
$+0.0001$ & $0.045629$ & $1.085485$ & $0.035\%$ & marginal \\
$+0.001$ & $0.046529$ & $1.082039$ & $0.353\%$ & excluded \\
$+0.005$ & $0.050529$ & $1.068400$ & $1.609\%$ & excluded \\
$-0.001$ & $0.044529$ & $1.088988$ & $0.364\%$ & excluded \\
\bottomrule
\end{tabular}
\end{table}

% ══════════════════════════════════════════════════════════════════════════════
\section{Epistemic status}
\label{sec:epistemic}

\begin{table}[ht]
\centering
\caption{Epistemic status of the results of D39.}
\label{tab:epistemic}
\medskip
\begin{tabular}{>{\raggedright}p{6.8cm}p{3.0cm}p{3.2cm}}
\toprule
Statement & Status & Source \\
\midrule
$\kap = P_1\cdot P_2$ (two-factor decomposition) &
  Algebraic identity & D39 Prop.~\ref{prop:decomp} \\[3pt]
$P_1 = \vp/3$ (active fraction of valence) &
  Theorem (D05) & C4 $\Rightarrow$ D05 \\[3pt]
$P(\mathrm{stable}|\mathrm{sea}) = 0$ &
  Theorem (D29) & $(A)\wedge(B)$ \\[3pt]
C3 alone $\Rightarrow$ uniform measure &
  \textbf{False} (counterexample) & D39 Prop.~\ref{prop:negative} \\[3pt]
H3 consistent with C1--C4 &
  Proved & D39 Prop.~\ref{prop:consistent} \\[3pt]
H3 unique under $S_4$-invariance &
  Proved & D39 Prop.~\ref{prop:unique} + D36 \\[3pt]
$\kap = 310\vp/11017$ under H3 &
  Theorem under H3 & D39 Theorem~\ref{thm:kappa} \\[3pt]
H3 from C1--C4 without D36 &
  \textbf{Open} (OP1 D39) & See below \\[3pt]
H3 constrained to $\pm 0.22\%$ &
  Observational bound & D36 sensitivity \\
\bottomrule
\end{tabular}
\end{table}

\begin{openprob}[OP1 D39 --- Axiom-level proof of H3]
\label{op:H3}
Derive H3 (uniform measure on $\Rtot$ from the perspective of an
external $\Kfour$ observer) from axioms C1--C4 alone, without
invoking the $S_4$-symmetry result of D36 as an additional input.

Concretely: show that the four \PDL\ axioms, applied to a proton
closure of type $(n_u, n_d, \rval, \Rsea, \Rtot)$, force the
cross-edge measure seen by $\Kfour$ to be uniform on $\Rtot$.

This would complete the derivation of $\kap$ from first principles
and close OP1 of D36 unconditionally.
\end{openprob}

% ══════════════════════════════════════════════════════════════════════════════
\section*{Conclusion}

This paper has established a systematic and complete analysis of
Open Problem~OP1 from D36.

\emph{The two-factor decomposition} $\kap = (\vp/3)\cdot(\rval/\Rtot)$
holds algebraically with machine-precision error. It separates OP1~D36
into two distinct sub-problems: the conditional engagement probability
$P_1 = \vp/3$ (resolved by D05) and the prior probability
$P_2 = \rval/\Rtot$ (requiring a measure on $\Rtot$).

\emph{The negative result} establishes that C3 alone does not supply
the required uniform measure: a constructive counterexample shows
that C3 is compatible with asymmetric measures.

\emph{Hypothesis~H3} is formulated as the Indifference Lemma: the
measure on $\Rtot$ is uniform because (i)~C3 makes all relations
candidates for engagement; (ii)~C4 makes the valence/sea partition
internal and invisible to $\Kfour$ before filtration; (iii)~the
$S_4$-symmetry of $\Kfour$ (D36) selects the uniform measure as
the unique $S_4$-invariant one.

\emph{Under H3}, $\kap = 310\vp/11017 \in \mathbb{Q}(\sqrt{5})$
is derived by the law of total probability from D05 and D29.
H3 is shown to be consistent with C1--C4 and constrained to
$\pm 0.22\%$ by the Hubble tension observable.

\emph{OP1 D36 is partially resolved}: the derivation is complete
under H3, and H3 is the unique irreducible residual. Its
unconditional proof from C1--C4 is Open Problem~OP1 of D39.

% ── Bibliography ─────────────────────────────────────────────────────────────
\begin{thebibliography}{99}

\bibitem{PDL_Main_2026}
C.~Laubscher,
\textit{The Projective Dynamic Logo: A Structural Derivation of
Fundamental Constants from Finite Signed Graphs},
Zenodo, 2026. \url{https://doi.org/10.5281/zenodo.14999037}

\bibitem{Combinatorial_Proton_PDL_2026}
C.~Laubscher,
\textit{Combinatorial Proton Architecture in the PDL Framework} (D05),
Zenodo, 2026.

\bibitem{PDL_D29_2026}
C.~Laubscher,
\textit{Axiomatic Derivation of the Leading Valence Engagement Amplitude
in the PDL Framework} (D29),
Zenodo, 2026. \url{https://doi.org/10.5281/zenodo.19283107}

\bibitem{PDL_D31_2026}
C.~Laubscher,
\textit{Proof of the Density-Dependent Gravitational Coupling
$G_{\mathrm{eff}}(N)=\sigma(N)\cdot G_{\mathrm{PDL}}$ in the PDL Framework}
(D31), Zenodo, 2026. \url{https://doi.org/10.5281/zenodo.19294984}

\bibitem{PDL_D36_2026}
C.~Laubscher,
\textit{Proof of $G_{\mathrm{eff}}(N)=\sigma(N)\cdot G_{\mathrm{PDL}}$
from the Trace Structure of $K_4$ and Independence of Gravitational
Engagement} (D36), Zenodo, 2026.
\url{https://doi.org/10.5281/zenodo.19323033}

\bibitem{PDL_D37_2026}
C.~Laubscher,
\textit{Surface Locality of Relational Entropy and the PDL Area Law}
(D37), Zenodo, 2026.

\bibitem{PDL_D38_2026}
C.~Laubscher,
\textit{Bekenstein--Hawking Entropy from the PDL Effective Gravitational
Coupling} (D38), Zenodo, 2026.

\end{thebibliography}

\end{document}